\documentclass[UTF8,a4paper,11pt,fontset=fandol]{ctexart}
\usepackage[margin=2.15cm]{geometry}
\usepackage{booktabs,tabularx,array,amsmath,graphicx,caption,fancyhdr,enumitem,xcolor,hyperref}
\hypersetup{hidelinks,pdfauthor={钱森}}
\pagestyle{fancy}\fancyhf{}\fancyfoot[R]{\thepage}\renewcommand{\headrulewidth}{0pt}
\setlength{\parskip}{3pt}\setlength{\emergencystretch}{2em}
\setlist{nosep,leftmargin=2em}
\newcommand{\titleline}[1]{\begin{center}{\LARGE\bfseries #1}\par\vspace{5pt}钱森\end{center}}
\newcommand{\fn}[1]{\nolinkurl{#1}}
\newcolumntype{L}{>{\raggedright\arraybackslash}X}


\begin{document}
\titleline{青河流域洪水预报系统需求规格说明}
\section{业务目标和使用者}
青河水文站的值班员接收六站雨量、蒸发皿蒸发量和出口观测流量，选定参数版本后起报。复核人员查看过程线、峰值和警戒状态，确认后签发产品。事后复盘需要回答三个问题：使用了哪一版输入，哪组参数产生了结果，谁在何时发起或签发了产品。课程系统在一个受信任的本地工作台上运行，操作员名称写入记录。
\section{五个核心用例}
\noindent\begin{tabularx}{\linewidth}{@{}p{2.5cm}LL@{}}\toprule 用例 & 主流程 & 可验收结果\\\midrule
UC1 入库与质控 & 导入示例或自备 CSV，提供初始蓄水与初始出流。逐行检查站点强迫、观测和时刻。 & 能按场次查到入库结果。问题标注行、字段及原因。阻断性错误禁止起报。\\
UC2 启动预报 & 选场次、方案及选项，立即运行或登记定时任务。 & 返回唯一预报编号、起报版本与逐时段流量；调度到期产生草稿。\\
UC3 结果与签发 & 打开历史预报，对比观测与基线，复核洪峰、预警及情景带，显式签发。 & 可查询产品、识别签发状态，并从编号恢复完整过程。\\
UC4 方案管理 & 新建参数集，对同名方案修订，切换活动版本，删除停用版本。 & 新版本有独立标识，后续起报引用所选版本，既有产品可追溯。\\
UC5 审计与回放 & 从历史记录读取输入快照和方案快照，重新执行并核对结果。 & 可查看回放结论与操作记录；历史产品与后来修改的活动参数相互独立。\\\bottomrule
\end{tabularx}
\section{四项扩展及验收边界}
\begin{description}
\item[O1 滚动校正] 每时段先输出先验预报，再以当时可用流量更新河道出流状态供下一步使用。截止时间之后的观测和缺测值不参与状态更新。修改未来观测不应改变较早时段的预测。
\item[O2 预警通知] 采用教学站两个流量阈值，达到500或900\,m³/s分别判为警戒或保证等级。超阈值起报应在站内消息中心留下关联预报的通知。
\item[O3 集合预报] 第二场使用五个降雨倍率与三个时移组合，共15个成员。并行计算后显示逐时段第5、50、95百分位。每个时段满足下界不大于中位数、中位数不大于上界。
\item[O4 Web 工作台] 可在浏览器完成导入、方案选择、起报、比较、签发与回放，直接进入产品详情，并查看质控问题和预警。
\end{description}
\section{数据和运行约束}
流域面积2000\,km²，步长6\,h，湿润地区教学洪水共四场。原始场次 CSV 的时间为流域当地时刻。六站加权得到面雨量，蒸发与降雨使用每时段毫米，出口流量使用 m³/s。缺失或负的必要强迫数据阻断起报，异常出口观测作为质控标记保存。独立检验采用第四场完整给定强迫。黄金测试逐个核验四场共160个流量结果。
\par 系统交付包括 Compose 启动配置、版本控制记录、CI、接口契约、黄金数据、测试与运维文档。数据卷应支持容器重建后查询原产品，备份还原后应能执行历史回放。
\section{关键数据与失败处理}
\noindent\begin{tabularx}{\linewidth}{@{}p{2.4cm}LL@{}}\toprule 对象 & 必需字段或条件 & 不满足条件时的行为\\\midrule
场次强迫 & 时间、青岭等六站雨量及蒸发皿读数。时刻严格按6小时递增，强迫为有限非负数。 & 结构性缺列返回校验错误；可定位的行级问题入库，问题场次禁止起报。\\
流量观测 & 观测列应存在，单元格可缺测。有效观测用于过程对比、精度计算和可选校正。 & 缺测标为警告；该点跳过观测更新与精度计算。\\
初始条件 & 上、下、深层土壤状态按 WU、WL、WD 提供，另需 S、QS、QI、QG。 & 七个值需有限且非负；状态超过当前方案容量时返回起报错误。\\
方案 & 16项完整参数及名称，传播参数满足 $2KE X\le6\le2KE(1-X)$。 & 越界参数拒绝保存，活动版本需在切换后才能归档。\\
签发产品 & 一次预报编号唯一对应一份结果，签发时记录署名、时间。 & 查询不存在的编号返回404。草稿下载公报返回409。\\\bottomrule\end{tabularx}
\section{现场验收场景}
先将缺测观测导入新场次，确认浏览器能定位问题并生成基准预报。随后以第二场为输入，运行15个情景并检查消息中心。复核人签发后下载公报，核对峰值、单位、时间轴和签发者。最后修改方案或更新观测，在旧编号上执行回放，核对结果摘要；再重启服务，用同一编号重复查询。验收以界面、HTTP响应和数据库持久化结果相互核对。

\end{document}
